\documentclass[11pt]{article}

\usepackage[margin=1in]{geometry}
\usepackage{amsmath, amssymb, amsthm, mathtools}
\usepackage{enumitem}
\usepackage{hyperref}

% Theorem environments
\newtheorem{theorem}{Theorem}[section]
\newtheorem{lemma}[theorem]{Lemma}
\newtheorem{proposition}[theorem]{Proposition}
\newtheorem{corollary}[theorem]{Corollary}

\theoremstyle{definition}
\newtheorem{definition}[theorem]{Definition}
\newtheorem{example}[theorem]{Example}

\theoremstyle{remark}
\newtheorem{remark}[theorem]{Remark}

% Document information
\title{Math 113: Lecture Notes}
\author{}
\date{November 17, 2026}

\begin{document}

\maketitle

\section{Lecture Notes}

\begin{theorem}[Kronecker's Theorem]
Let $F$ be a field and $f(x) \in F[x]$ a nonconstant polynomial.
Then there exists an extension field $E \supseteq F$
and $\alpha \in E$ such that $f(\alpha) = 0$.
\end{theorem}

\begin{proof}
We know that f(x) can be factored into irreducible polynomials over F. Pick p(x) to be on eo fhtese irreducible factors. It suffices to find a zero $\alpha$ of p(x), since this will also be a zero of f(x)
\\From last time, we know $\langle p(x) \rangle$ is maximal, Since E is a field, and $\langle p(x) \rangle$ is maximal. 
\\ Want injective ring homomorphism $\psi:F \rightarrow E = F[x] / \langle p(x) \rangle$
\\ Let's take $a \mapsto a + \langle p(x) \rangle$
\\This is a ring homomorphism, since it is the composition of homomorphisms. $F \rightarrow F[X] \rightarrow F[x] / \langle p(x) \rangle$
\\However, we need to show that this is injective. 
\\Suppose $\psi(a) = \psi(b), \text{then } a+\langle p(x) \rangle = b _ \langle p(x) \rangle$, so $a-b \in \langle p(x) \rangle$
\\a-b is constant, so a-b = 0, so a = b and $\psi $ is injective
\\ So we identity F with $\operatorname{im}{\psi} \subseteq E$


Let $\alpha = x + \langle p(x) \rangle$ 

$\phi_{\alpha}(p(x)) = 0$

Let $p(x) = a_0 + a_1x + \dots + a_nx^n, a_i \in F \subseteq E$
$phi_{\alpha}(p(x)) = a_0 + a_1(x + \langle p(x) \rangle) + \dots + a_n(x+ \langle p(x) \rangle)^n$
$= (a_0 + a_1x + a_1x^2 + \dots + a_nx^n) + \langle p(x) \rangle = p(x) + \langle p(x) \rangle = 0 + \langle p(x) \rangle = 0 \in E$. Thus we found $\alpha \in E \text{ with } p(\alpha) = 0$

\end{proof}

\begin{example}
    $F = \mathbb{R}, f(x) = x^2 + 1$ Note that this polynomial is irreducible over the field. 
    
    $E = \mathbb{R}[x] / \langle x^2 + 1 \rangle \text{ is a field. And } x + \langle x^2 + 1 \rangle \text{ is a zero of f}$

    $\phi: E \cong \mathbb{C}, \phi(x + \langle x^2 + 1 \rangle) = i$

    $i^2 = x^2 + \langle x^2 + 1 \rangle \Rightarrow i^2 + 1 = x^2 + 1 + \langle x^2 + 1 \rangle = 0$

\end{example}

\begin{example}
    $F = \mathbb{Q}, f(x) = x^4 - 5x^2 + 6 = (x^2-2)(x^2-3)$
    
    $p(x) = x^2 -2$ or $p(x) = x^2 - 3$

    
\end{example}
\end{document}

